% !TeX root = ../../Book1.tex
% 纲要标题：### 9.2 挠与无挠
% 纲要位置：计划纲要.md，第 364 行。

\section{挠与无挠}
\label{b1:sec:0902}

% 纲要内容（仅作写作计划，尚非教材正文）：
%
% * 挠子群与商群
% * 有限阿贝尔群的 p 主分解
% * 有限生成条件不可省略的例子
%

\subsection{挠子群与商群}
\label{b1:subsec:0902:01}

整数格中非零向量没有有限阶，但一般阿贝尔群可以同时含有无限阶和有限阶元素。先把后一部分提取出来，可以把后续分类拆成自由方向与有限阶方向。

\begin{definition}\label{b1:def:torsion-subgroup}
阿贝尔群 \(A\) 中满足某个正整数 \(m\) 使 \(mx=0\) 的元素称为\term[naoyuan]{挠元}[torsion element]。全部挠元组成的集合记为
\[
A_{\mathrm{tor}}\coloneqq\{\,x\in A\mid mx=0
\text{ 对某个整数 }m\ge1\,\}.
\]
若 \(A_{\mathrm{tor}}=A\)，称为\term[naoqun]{挠群}[torsion group]；若 \(A_{\mathrm{tor}}=0\)，称为\term[wunao-qun]{无挠群}[torsion-free group]。
\end{definition}
\symidx{\(A_{\mathrm{tor}}\)：挠子群}

\begin{proposition}\label{b1:prop:torsion-quotient}
\(A_{\mathrm{tor}}\) 是 \(A\) 的特征子群，称为\term[naoziqun]{挠子群}[torsion subgroup]；商群 \(A/A_{\mathrm{tor}}\) 无挠。
\end{proposition}
\begin{proof}
若 \(mx=0,ny=0\)，则交换性给出 \(mn(x-y)=0\)，故挠元对差封闭。同态保持整数倍，所以自同构保持挠元，给出特征性。若 \(k(x+A_{\mathrm{tor}})=0\)，则 \(kx\) 为挠元，存在 \(m\ge1\) 使 \(mkx=0\)。于是 \(x\in A_{\mathrm{tor}}\)，原陪集就是零陪集。
\end{proof}
交换性不可随意去掉。在无限二面体群中，两次反射都有二阶，而其乘积可以是无限阶平移。因此一般群的有限阶元素集合未必为子群。

\begin{proposition}\label{b1:prop:fg-abelian-subgroups}
有限生成阿贝尔群的每个子群都有限生成；其中挠子群为有限群。
\end{proposition}
\begin{proof}
取满同态 \(\pi:\mathbb Z^n\to A\)。对子群 \(B\le A\)，原像 \(\pi^{-1}(B)\le\mathbb Z^n\) 由前节定理有有限基，其像生成 \(B\)。特别地，\(A_{\mathrm{tor}}\) 有有限生成元 \(t_1,\ldots,t_s\)。若 \(t_i\) 的阶为 \(m_i\)，每个挠元都能写成
\(\sum a_i t_i\)，其中 \(0\le a_i<m_i\)，故元素数至多 \(\prod m_i\)，为有限群。
\end{proof}
在 \(A=\mathbb Z\oplus C_6\) 中，\((u,\bar v)\) 为挠元当且仅当 \(u=0\)，所以挠子群为 \(C_6\)，无挠商群为 \(\mathbb Z\)。一般有限生成阿贝尔群是否都能像这样拆成直和，仍需下一节的结构定理。

\subsection{有限阿贝尔群的主分解}
\label{b1:subsec:0902:02}

不同素数的挠可以独立拆开。对素数 \(p\)，定义
\[
A_p=\{\,x\in A\mid p^k x=0\text{ 对某个 }k\ge0\,\},
\]
称为 \(A\) 的\term[p-zhufenliang]{\(p\) 主分量}[p-primary component]。它是特征子群，因为两个元素可取同一个较大的 \(p\) 幂来消去，且同态保持其消去关系。
\symidx{\(A_p\)：阿贝尔群的 \(p\) 主分量}

\begin{theorem}\label{b1:thm:abelian-primary}
设 \(A\) 为有限阿贝尔群，\(|A|=p_1^{e_1}\cdots p_t^{e_t}\)。则加法映射给出同构
\[
A_{p_1}\oplus\cdots\oplus A_{p_t}\ \cong\ A.
\]
\end{theorem}
\begin{proof}
若 \(A=0\)，素因子列表为空，空直和为零群，结论成立。以下设 \(A\ne0\)。令 \(N=|A|\)、\(N_i=N/p_i^{e_i}\)。这些 \(N_i\) 的最大公因数为一，故 Bézout 等式给出整数 \(c_i\) 使 \(\sum_i c_iN_i=1\)。对任意 \(x\in A\)，Lagrange 定理给出 \(Nx=0\)，所以
\[
x=\sum_i x_i,\qquad x_i=c_iN_i x,\qquad
p_i^{e_i}x_i=0.
\]
这说明各主分量之和为 \(A\)。

为证唯一性，设 \(\sum_i y_i=0\)，其中 \(y_i\in A_{p_i}\)。每个 \(y_i\) 的阶既为 \(p_i\) 幂又整除 \(N\)，所以 \(p_i^{e_i}y_i=0\)。对等式施以整数倍 \(c_jN_j\)：当 \(i\ne j\) 时，\(N_j\) 含因子 \(p_i^{e_i}\)，故对应项为零；当 \(i=j\) 时，由 Bézout 等式模 \(p_j^{e_j}\) 得 \(c_jN_j\equiv1\)，对应项仍为 \(y_j\)。于是 \(y_j=0\)，对每个 \(j\) 都成立。
\end{proof}
这称为有限阿贝尔群的\term[zhufenjie]{主分解}[primary decomposition]。例如在 \(C_{12}\) 中，对 \(N_2=3,N_3=4\)，可取 \(c_2=-1,c_3=1\)。所以每个 \(\bar x\) 唯一写为
\[
\bar x=(-3\bar x)+(4\bar x),
\]
前项被 \(4\) 消去，后项被 \(3\) 消去，对应 \(C_{12}\cong C_4\oplus C_3\)。这种投影是由群上的整数倍映射给出的，不需要预先选择每个主分量的生成元。

\subsection{有限生成条件与反例}
\label{b1:subsec:0902:03}

有限生成条件的作用至少有两处：它让关系由有限矩阵描述，也让挠子群有限。两个具体例子可以防止把有限生成结构定理误用于一般阿贝尔群。

\ProseParagraphBreak
首先，\((\mathbb Q,+)\) 无挠，但不有限生成。若有限多个有理数的分母都整除 \(N\)，它们的整数线性组合都在 \(\frac1N\mathbb Z\) 中；有理数 \(\frac1{2N}\) 不在其中。它也不是自由阿贝尔群：在 \(\mathbb Q\) 中每个元素都可除以二，故 \(\mathbb Q=2\mathbb Q\)；在任何非零自由阿贝尔群中，一个基元素都不属于二倍子群。这个区别排除了任意有限或无限秩的自由基。

\ProseParagraphBreak
其次，固定素数 \(p\)，考虑
\[
C_{p^\infty}
=\mathbb Z[1/p]/\mathbb Z,\qquad
\mathbb Z[1/p]=\{\,a/p^k\mid a\in\mathbb Z,\ k\ge0\,\}.
\]
这里 \(\mathbb Z[1/p]\) 只作为有理数加法的子群使用。商群中 \(\frac1{p^k}+\mathbb Z\) 的阶恰为 \(p^k\)，所以它是无限挠群。有限多个元素都在某一个由 \(\frac1{p^k}+\mathbb Z\) 生成的有限循环群中，因而整个群不有限生成。这个群称为\term[prufer-qun]{Prüfer 群}[Prüfer group]。

\ProseParagraphBreak
上述两个群还共享一种现象：有限生成的小部分很简单，全部小部分的并却未必具有同样的有限描述。因此后续每次使用“有限个不变因子”或“挠子群有限”时，都应保留有限生成这一条件。下一节证明的无挠推出自由，也明确限于有限生成阿贝尔群。
